% ---
% §1.5 Alcance y limitaciones. Reemplaza las líneas 662–688 de Docs/Technical_Report.tex.
% Cambios respecto a TT1:
%   - propósito y pregunta alineados con la pregunta reformulada (§1.2, D-041);
%   - se añaden las codificaciones evaluadas y excluidas (decidido con el autor el 2026-10-15)
%     y las limitaciones de datos y de diseño;
%   - se quitan la analogía del ENIAC y el párrafo de Grover (auditoría T3: fuera de tema y sin
%     cita; la afirmación sobre NISQ y Grover no tenía fuente). El texto original sigue en el .tex.
% Decisiones y hallazgos que lo respaldan: D-005, D-013, D-014, D-031; H-002, H-035, H-045, H-051,
% H-052, H-054.
% ---

\subsection{Alcance y limitaciones}
\label{sec:alcance}

El presente trabajo tiene como propósito central evaluar la viabilidad de las transformaciones
cuánticas variacionales como mecanismo de codificación de características radiómicas extraídas de
imágenes mamográficas, determinando en qué medida modifican la separabilidad de clases
benigno/maligno frente a representaciones clásicas construidas a partir del mismo vector de
características. El sistema desarrollado constituye un prototipo de investigación académica, no un
sistema de diagnóstico clínico, y su alcance está delimitado por los criterios siguientes.

\textbf{Datos y tarea}

El sistema opera exclusivamente sobre las lesiones del dataset CBIS-DDSM que corresponden a masas y
calcificaciones con diagnóstico patológico confirmado, a partir de la región de interés que delimita
la máscara de cada lesión. Masas y calcificaciones se tratan como subproblemas independientes en
todo el \textit{pipeline}. El problema abordado es la clasificación binaria benigno/maligno; la
detección de lesiones, su localización espacial y la segmentación quedan fuera del alcance.

\textbf{Simulación clásica}

% respaldo: «“Noisy” emphasizes that we’ll have imperfect control over those qubits; the noise will place serious limitations on what quantum devices can achieve in the near term.» (Preskill2018NISQ, p. 5)
La totalidad de los experimentos cuánticos se ejecuta en simulación clásica con Qiskit Aer, sin
acceso a \textit{hardware} cuántico. Los resultados corresponden, por tanto, a condiciones ideales,
libres del ruido de los dispositivos físicos, lo cual es una condición deliberada del diseño
experimental y no una deficiencia metodológica: permite atribuir lo observado a la codificación y no
al ruido. Los dispositivos actuales, de escala intermedia y con ruido (NISQ), tienen un control
imperfecto de sus qubits, y ese ruido limita lo que pueden lograr a corto plazo
\cite{Preskill2018NISQ}. Lo que este trabajo dice sobre la ejecución en un procesador real se basa en
proyecciones: el número de mediciones se estudia con un muestreo ideal de las probabilidades
exactas, sin modelar el ruido del dispositivo.

\textbf{Codificaciones evaluadas}

El estudio se limita a dos mapas de características con entrelazamiento entre todos los pares de
qubits, \textit{ZZFeatureMap} y \textit{PauliFeatureMap} con los términos Y y ZZ, con una repetición
y $k = 12$ qubits. El número de qubits se fijó por el costo de simulación y por la diversidad de
familias de características que conserva la selección, y obliga a descartar parte de la información
de las 67 características extraídas. Quedan fuera las codificaciones sin entrelazamiento, como la
angular producto o la codificación en amplitudes, porque su kernel es por construcción un kernel
clásico (sección~\ref{sec:modulo5}). Los parámetros de la capa variacional se fijan con una semilla
y no se entrenan, para que la separabilidad se mida sobre una representación no supervisada;
entrenarlos se explora solo como experimento complementario.

\textbf{Datos de imagen}

Las mamografías del CBIS-DDSM son películas digitalizadas sin escala física de intensidad ni
espaciado de píxel en sus metadatos, de modo que las características de primer orden se calculan
sobre intensidades no calibradas y las de forma se expresan en píxeles. Todos los resultados
provienen de un único conjunto de datos, sin validación externa, y el conjunto de prueba oficial se
usa una sola vez por modelo.

\textbf{El objetivo no es superar lo existente, sino explorar lo posible}

Un error frecuente en la evaluación de trabajos de aprendizaje automático cuántico aplicado consiste
en juzgar su validez en función de si el modelo cuántico supera en precisión a los mejores modelos
clásicos disponibles. Este criterio, pertinente para tecnologías maduras, resulta inapropiado para el
estado actual de la computación cuántica. El objetivo del presente trabajo no es demostrar que las
representaciones cuánticas superan a una CNN de última generación en AUC-ROC, sino responder la
pregunta planteada en la sección~\ref{sec:planteamiento}: en qué medida estas transformaciones
modifican la separabilidad intrínseca de las clases sobre datos mamográficos reales, y si esa
diferencia se refleja en la clasificación. Una mejora en separabilidad sin mejora en la clasificación
constituiría evidencia válida sobre el comportamiento de estos circuitos, y también lo es una
separabilidad igual o menor que la clásica, siempre que esté medida con el mismo rigor y que se
documenten las restricciones que la explican.
